\section{Introduction}\label{intro}

Let $\F$ be any field and $X=\{x_1,x_2,\ldots,x_n\}$ be a set of $n$
free noncommuting variables. Let $X^*$ denote the set of all free
words (which are monomials) over the alphabet $X$ with concatenation
of words as the monoid operation and the empty word $\epsilon$ as
identity element.

The \emph{free noncommutative ring} $\FX$ consists of all finite
$\F$-linear combinations of monomials in $X^*$, where the ring
addition $+$ is coefficient-wise addition and the ring multiplication
$*$ is the usual convolution product. More precisely, let $f,g\in\FX$
and let $f(m)\in\F$ denote the coefficient of monomial $m$ in
polynomial $f$. Then we can write $f=\sum_m f(m) m$ and $g=\sum_m g(m)
m$, and in the product polynomial $fg$ for each monomial $m$ we have

\[
fg(m)=\sum_{m_1m_2=m} f(m_1)g(m_2).
\]

The \emph{degree} of a monomial $m\in X^*$ is the length of the
monomial $m$, and the degree $\deg f$ of a polynomial $f\in\FX$ is the
degree of a largest degree monomial in $f$ with nonzero coefficient.
For polynomials $f,g\in\FX$ we clearly have $\deg (fg) = \deg f + \deg
g$.

A \emph{nontrivial factorization} of a polynomial $f\in\FX$ is an
expression of $f$ as a product $f=gh$ of polynomials $g,h\in\FX$ such
that $\deg g > 0$ and $\deg h > 0$. A polynomial $f\in\FX$ is
\emph{irreducible} if it has no nontrivial factorization and is
\emph{reducible} otherwise.  For instance, all degree $1$ polynomials
in $\FX$ are irreducible. Clearly, by repeated factorization every
polynomial in $\FX$ can be expressed as a product of
irreducibles. 

In this paper we study the algorithmic complexity of polynomial
factorization in the free ring $\FX$. The factorization algorithm is
by an application of Higman's linearization process followed by
factorization of a matrix with linear entries (under some technical
conditions) using Cohn's factorization theory.

It is interesting to note that Higman's linearization process
\cite{Hig} has been used to obtain a deterministic polynomial-time
algorithm for the RIT problem. That is, the problem of testing if a
noncommutative rational formula (which computes an element of the free
skew field $\skewf$) is zero on its domain of definition
\cite{GGOW20,IQS17, IQS18,HW15}.

%

\subsection{Overview of the results}

The main result of the paper is the following.

\begin{theorem*}[Main Theorem]
Given a multivariate noncommutative polynomial $f\in\F_q\angle{X}$ for
a finite field\footnote{We present the detailed randomized algorithm
  over large finite fields. In the case of small finite fields we
  obtain a deterministic $\poly(s,q,|X|)$ time algorithm with minor
  modifications.} $\F_q$ by a
noncommutative algebraic branching program of size $s$ as input, a
factorization of $f$ as a product $f=f_1f_2\cdots f_r$ can be computed
in randomized time $\poly(s,\log_2q,|X|)$, where each
$f_i\in\F_q\angle{X}$ is an irreducible polynomial that is output as
an algebraic branching program.
\end{theorem*} 
%


The proof has three broad steps described below.

\begin{itemize}
\item \textbf{Higman linearization and Cohn's factorization theory}~~
  Briefly, given a noncommutative polynomial $f\in\FX$, we can
  transform it into a linear matrix $L$ such that $f\oplus I = PLQ$,
  where $P$ is an upper triangular matrix with polynomial entries and
  all $1$'s diagonal and $Q$ is a lower triangular matrix with
  polynomial entries and all $1$'s diagonal, $P$ and $Q$ are the
  matrices implementing the sequence of row and column operations
  required for the Higman linearization process. Now, as observed
  by Garg et al \cite{GGOW20} (in their work on rational identity
  testing), if $f$ is given by a noncommutative formula as input
  then the above Higman linearization can be carried out in polynomial
  time. We strengthen their observation with a modified Higman
  linearization process that we call Block-Higman linearization (to
  emphasize that the modification works with matrix blocks) and
  show the more general result that if $f$ is given by an ABP as input
  we can still compute its Higman linearization as defined above
  in polynomial time.

%

\item \textbf{Ronyai's common invariant subspace algorithm}~~ Next,
  the most important tool algorithmically, is Ronyai's algorithm for
  computing common invariant subspaces of a collection of matrices
  over finite fields \cite{Ronyai2}. We show that Ronyai's common
  invariant subspace algorithm can be repeatedly applied to factorize
  a linear matrix $L=A_0+\sum_{i=1}^n A_i x_i$, into a product of
  irreducible linear matrices provided $A_0$ is invertible and
  $[A_1 A_2\cdots A_n]$ has full row rank or
  $[A_1^T A_2^T\cdots A_n^T]^T$ has full column rank. The later
  conditions are called as right and left monicity of the linear
  matrix $L$ respectively. With some technical work we can ensure
  these conditions for a linear matrix $L$ that is produced from a
  polynomial $f$ by Higman linearization. Then Ronyai's algorithm
  yields the factorization of $L$ into a product of irreducible linear
  matrices (upto multiplication by units). Here, Cohn's theory of
  factorization of noncommutative linear matrices gives us sufficient
  useful information about the structure of irreducible linear matrices.

\item \textbf{Recovering the factors of $f$}~~ Finally, we design a
  simple linear algebraic algorithm for trivializing a matrix product
  $AB=0$, where $A$ is a linear matrix and $B$ is a column vector of
  polynomials from $\FX$, using which we are able to extract the
  irreducible factors of $f$ from the factors of $L$. An invertible
  matrix $M$ with polynomial entries \emph{trivializes} the relation
  $AB=0$ if the modified relation $(AM)(M^{-1}B)=0$ has the property
  that for every index $i$ either the $i^{th}$ column of $AM$ is zero
  or the $i^{th}$ row of $M^{-1}B$ is zero. While such matrices $M$
  exist for any matrix product $AB=0$ with entries from $\FX$, we
  obtain an efficient algorithm in the special case when $A$ is linear
  and $B$'s entries are polynomials computed by small arithmetic
  circuits. This special case is sufficient for our application.
\end{itemize}

There are some additional technical aspects we need to deal with. 
%
%
%
%
Let $L=A_0+\sum_{i=1}^n A_i x_i$ be the linear matrix obtained from
$f\in\F_q\angle{X}$ by Higman linearization, where
$X=\{x_1,x_2,\ldots,x_n\}$ and $A_i\in \F_q^{d\times d}, 0\le i\le
n$. If $A_0$ is an invertible matrix then it turns out that the
problem of factorizing $L$ can be directly reduced to the problem of
finding a common invariant subspace for the matrices $A_0^{-1}A_i,
1\le i\le n$.  In general, however, $A_0$ is not invertible. Two cases
arise:
\begin{itemize}
\item[(a)] The polynomial $f$ is \emph{commutatively nonzero}. That
  is, it is nonzero on $\F_q^n$ (or on $\F^n$ for a small extension
  field $\F$). In this case, by the DeMillo-Lipton-Schwartz-Zippel
  Lemma~\cite{DL78,Sch80, Zip79}, we can do a linear shift of the
  variables $x_i\leftarrow x_i+\alpha_i$ in the polynomial $f$, for
  $\alpha_i$ randomly picked from $\F_q$ (or $\F$). Let the resulting
  polynomial be $f'$ and let its Higman linearization be $L_{f'}$. In
  $L_{f'}$ the constant matrix term $A'_0$ will be invertible with
  high probabilty, and the reduction steps outlined above will work
  for $L_{f'}$. Furthermore, from the factorization of $f'$ we can
  efficiently recover the factorization of $f$. Section~\ref{cnz-sec}
  deals with Case~(a), with Theorem~\ref{cnzthm} summarizing the
  algorithm for factorizing $f$. Theorem~\ref{lfact1} describes the
  algorithm for factorization of the linear matrix $L_{f'}$, and the
  factor extraction lemma (Lemma~\ref{extract}) allows us to
  efficiently recover the factorization of $f'$ from the factorization
  of $L_{f'}$.
  
\item[(b)] In the second case, suppose $f$ is zero on all scalars. Then, for example by
  Amitsur's theorem \cite{Ami66}, for a random matrix substitution
  $x_i\leftarrow M_i\in \F^{2s\times 2s}$ the matrix
  $f(M_1,M_2,\ldots,M_n)$ is \emph{invertible} with high probability,
  where $s$ is the formula size of $f$.\footnote{Amitsur's theorem
    strengthens the Amitsur-Levitski theorem \cite{AL50} often used in
    noncommutative PIT algorithms \cite{BW05}.} \footnote{In the actual algorithm we pick the matrices $M_i$ using a result from  \cite{DM17}}   Accordingly, we can
  consider the factorization problem for shifted and dilated linear
  matrix $L'=A_0\otimes I_{\ell} + \sum_{i=1}^n A_i\otimes (Y_i+M_i)$ which             
  will have the constant matrix term invertible, where each $Y_i$ is  
  an $\ell \times \ell$ matrix of distinct noncommuting variables, where $\ell=2s$.  Recovering
  the factorization of $L$ from the factorization of $L'$ requires
  some additional algorithmic work based on linear algebra. A lemma
  from \cite{HKV20} (refer Section \ref{cz-sec} and the Appendix for the details) turns
  out to be crucial here. The algorithm handling Case~(b) is described
  in Section~\ref{cz-sec}. Indeed, the new aspect of the algorithm is
  factorization of the dilated matrix $L'$ from which we recover the
  factorization of the Higman linearization $L_f$ of $f$. The
  remaining algorithm steps are exactly as in Section~\ref{cnz-sec}.
\end{itemize}  

\subsection{Small Finite fields}\label{small-field}

We now briefly explain the deterministic $\poly(s,q,|X|)$ time
factorization algorithm (when $\F_q$ is small). There are two places
in the factorization algorithm outlined above where randomization is
used: first, to obtain a matrix tuple $(M_1,M_2,\ldots,M_n)$ such that
$f(M_1,M_2,\ldots,M_n)$ is invertible, which ensures that the constant
matrix term of the linear matrix $L'$ is invertible. When
$q=\Omega(d)$, where $d= \deg f$, it suffices to randomly pick
$M_i\in\F_q^{2s\times 2s}$. However, if $q<d$ we can choose entries of
the matrices $M_i$ from a small extension field $\F_{q^k}$ such that
$q^k=\Omega(d)$.  Thereby, we will obtain factorization of $L'$ and
subsequently that of the polynomial $f$ over the extension field
$\F_{q^k}$. However, we can use the fact that the finite field
$\F_{q^k}$ can be embedded using the regular representation of the
elements of $\F_{q^k}$ in the matrix algebra $\F_q^{k\times k}$. Thus,
we can obtain from $(M_1,M_2,\ldots,M_n)$ a matrix tuple
$(M'_1,M'_2,\ldots,M'_n)$ with $M'_i\in\F_q^{2sk\times 2sk}$ such that
$f(M'_1,M'_2,\ldots,M'_n)$ is invertible. This will ensure that the
linear matrix $L'$ can be factorized over the field $\F_q$ which will
allow us to obtain a complete factorization of $f$ into irreducible
factors over $\F_q$.

In order to get a deterministic polynomial-time algorithm for finding
such matrices $M'_i, 1\le i\le n$ we will use the fact that the
polynomial $f$ is given by a small noncommutative formula and hence
has a small algebraic branching program. Then, using ideas from
\cite{RS05,F14,ACDM20} we can easily find such matrices $M'_i$ in
deterministic polynomial time.

Next, we notice that Ronyai's algorithm for finding common invariant
subspaces of matrices over $\F_q$ is essentially a polynomial-time
reduction to univariate polynomial factorization over $\F_q$. We can
use Berlekamp's deterministic $\poly(q,D)$ algorithm for the
factorization of univariate degree $D$ polynomials over
$\F_q$. Putting it together, we can obtain a deterministic
$\poly(s,q,|X|)$ time algorithm for factorization of
$f\in\F_q\angle{X}$ as a product of irreducible factors over
$\F_q$.\smallskip


\begin{comment}
Infact, it turns out that, building on ideas from Arvind
et.al. \cite{ACDM20} combined with determinsitic algorithm for PIT of
Read Once ABPs by Forbes and Shpilka \cite{} and using regular
representation of $\F_{q^k}$ as $k \times k$ matrices over
$\mathbb{F}_q$ we can efficiently find a matrix substitution
$\overline{N}= (N_1, N_2, \ldots, N_n)$ such that entries of $N_i$'s
are from base field $\F_q$ and the size of $N_i$ is at most $k$ times
as that of the original substitution matrices $M_i$'s and the linear
matrix obtained by using the substitution $\overline{N}$ instead of
$\overline{M}$ also has invertible constant term. As a result we can
obtain factorization of $f$ over $\F_q$. Moreover, for the case of
small finite fields, Ronyai's algorithm works in deterministic
polynomial time. Putting it together, in the case of the small finite
fields, we get deterministic polynomial time factorization algorithm
for factorization of the polynomial $f$ given by noncommutative
formula. 
\end{comment}

%

\subsection{Finite fields versus Rationals}\label{finite-vs-rationals}

Unfortunately, the algorithm outlined above does not yield an
efficient algorithm for noncommutative polynomial factorization over
rationals. The bottlneck is the problem of computing common invariant
subspaces for a collection of matrices over $\Q$. Ronyai's algorithm
for the problem over finite fields \cite{Ronyai2} builds on the
decomposition of finite-dimensional associative algebras over
fields. Given an algebra $\mathcal{A}$ over a finite field $\F_q$ the
algorithm decomposes $\mathcal{A}$ as a direct sum of minimal left
ideals of $\mathcal{A}$ which is used to find nontrivial common
invariant subspaces. However, as shown by Friedl and Ronyai
\cite{FR85}, over rationals the problem of decomposing a \emph{simple}
algebra as a direct sum of minimal left ideals is at least as hard as
factoring square-free integers.


\subsection{Related research}

The study of factorization in noncommutative rings is systematically
investigated as part of Cohn's general theory of noncommutative free
ideal rings \cite{Cohnfir,Cohnintro} which is based on the notion of
the weak algorithm. In fact, there is a hierachy of weak algorithms
generalizing the division algorithm for commutative integral domains
\cite{Cohnfir}.\smallskip

\noindent\textit{Algorithmic:}~~To the best of our knowledge, the
complexity of noncommutative polynomial factorization has not been
studied much, unlike the problem of commutative polynomial
factorization \cite{vzg-book,K89,KT90}.  Prior work on the complexity
of noncommutative polynomial factorization we are aware of is
\cite{AJR18} where efficient algorithms are described for the problem
of factoring \emph{homogeneous} noncommutative polynomials (which
enjoy the unique factorization property, and indeed the algorithms in \cite{AJR18} crucially use the unique factorization property). When the input homogeneous
noncommutative polynomial has a small noncommutative arithmetic
circuit (even given by a black-box as in Kaltofen's algorithms
\cite{K89, KT90}) it turns out that the problem is efficiently
reducible to commutative factorization by set-multilinearizing the
given noncommutative polynomial with new commuting variables. This
also works in the black-box setting and yields a randomized
polynomial-time algorithm which will produce as output black-boxes for
the irreducible factors (which will all be homogeneous). When the
input homogeneous polynomial is given by an algebraic branching
program there is even a deterministic polynomial-time factorization
algorithm. Indeed, the noncommutative factorization problem in for
homogeneous polynomials efficiently reduces to the noncommutative PIT
problem \cite{AJR18}, analogous to the commutative case \cite{KSS14},
modulo the randomness required for univariate polynomial factorization
in the case of finite fields of large characteristic. The motivation
of the present paper is to extend the above results to the
inhomogeneous case.\smallskip

\noindent\textit{Mathematical:}~~From a mathematical perspective,
building on Cohn's work there is a lot of research on the study of
noncommutative factorization. For example, \cite{BS15,BHL17} focus on
the lack of unique factorization in noncommutative rings and study the
structure of multiple factorizations. The research most relevant to
our work is the study of noncommutative analogues of the
Nullstellensatz by Helton, Klep and Volcic \cite{HKV18,HKV20}. In  
these papers the authors study the free singularity locus of a
noncommutative polynomial $f\in\FX$ where $\F$ is an algebraically
closed field of characteristic zero (in \cite{HKV20} mostly they
consider complex numbers). This is the set of all matrix tuples
$\bar{M}\in\mathcal{L}_n(f)$ (in all matrix dimensions $d$) where
$\mathcal{L}_n(f)=\{\bar{M}\mid \det f(\bar{M}) = 0,$ where $\bar{M}$
is an $n$-tuple of matrices$\}$.  It turns out that $f\in\FX$ is
irreducible if and only if for all $d\ge d_0$ for some $d_0$ the
hypersurface $\mathcal{L}_d(f)$ is irreducible which in turn holds iff
$\det f(\bar{X})$ is an irreducible commutative polynomial, where
$\bar{X}$ are generic matrices with commuting variables of dimension
$d\ge d_0$. However, $d_0$ turns out to be exponentially large.



\smallskip

\noindent\textbf{Plan of the paper.}~ In Section~\ref{prelim} we
present basic definitions and the background results from Cohn's work on factorization. In Section~\ref{bas-res-sec} we further present some results from Cohn's work relevant to the paper. In
Section~\ref{cnz-sec} we present the factorization algorithm for
polynomials $f$ that does not vanish on scalars and in
Section~\ref{cz-sec} we present the algorithm for the general case.
